\answer{ex:05:1}{$\ell\approx17\um$、$\approx100$日。脳全体への配信を担うのは配線の分岐であり、拡散は各放出部位を$\sim20\um$ぼかすだけである。}
\answer{ex:05:2}{$r=ab/(1+G)$、$S=1/(1+G)$は厳密。$+20\,\%$ではなく$+1.7\,\%$。}
\answer{ex:05:3}{$T^*$は$G=2$で$1.21\,\text{s}$、$G=9$で$0.19\,\text{s}$。現実的には$G\sim2$--$4$で、抑制は$3$--$5$分の1。}
\answer{ex:05:4}{総レートを$\lambda$として$\mathrm{CV}=1/\sqrt{2\lambda\tau_c}\approx9\cdot10^{-4}$。ニューロン1個なら$\tau_c=12.5\,\text{s}$が必要。}
\answer{ex:05:5}{どちらの場合も$c\approx0.40$。$\mathrm{CV}=0.050$対$0.158$で、$\sqrt{10}$分の1。個々のニューロンのレートは$a_i$に比例したままで、調節されるのは和だけである。}
\answer{ex:05:6}{$N_1=\overline{A\vee B}$、$N_2=\overline{A\vee N_1}$、$N_3=\overline{B\vee N_1}$、$N_4=\overline{N_2\vee N_3}$、XOR $=N_5=\overline{N_4}$。各切り替えに$\tau_c\ln2$かかり、ゲート$4$段の経路でXOR 1回に$2.8\,\text{s}$。}
\answer{ex:05:7}{(a)~$\tau_\gamma=6/\ln3\approx5.5\,\text{s}$。(b)~$D<2.2\,\text{s}$に対して$\bar D<4\,\text{s}$。遅延が予測できないと忍耐強くなる。}
\answer{ex:05:8}{$k_2=1/\tau_d=10\,\text{s}^{-1}$、$k_1=1/\tau_d-1/\tau_\gamma=9.5\,\text{s}^{-1}$。$\tau_\gamma=1/(k_2-mk_1)$なので、$+2.6\,\%$で2倍になる（$+5.3\,\%$で時間幅は無限大）。}
